% Related work (owner SK). Only sources actually consulted and cited in refs.bib
% appear here; the scope limits of this survey are stated rather than hidden.
\section{Related Work}
\label{sec:related}

\subsection{Anti-spoofing benchmarks and their language}

ASVspoof 2019 LA~\cite{todisco2019asvspoof} is the reference benchmark for logical
access attacks and the training source for our Stage-1 system. It is English, read
in studio conditions, and its attacks are the text-to-speech and voice conversion
systems of its era. Its influence is what makes the question in this paper worth
asking: a field that tunes against one language in one recording condition can
mistake in-domain accuracy for detection ability.

\subsection{Self-supervised representations for detection}

The current recipe is to fine-tune a self-supervised speech encoder rather than
engineer features. wav2vec~2.0~\cite{baevski2020wav2vec} learns from unlabelled audio
and transfers well to spoofing detection, where using it as a front-end gives large
gains over hand-designed features~\cite{tak2022wav2vec}. We adopt that architecture
directly: a wav2vec~2.0 encoder with attentive statistics
pooling~\cite{okabe2018attentive} and a linear classifier.

What this literature does not settle is what the encoder's transfer depends on. Our
probing results (Section~\ref{sec:probe}) find the decision forming within the first
three layers, before linguistic representation exists --- which is consistent with a
detector reading low-level synthesis artefacts, and is why a change of channel hurts
it more than a change of language.

\subsection{Parameter-efficient adaptation}

LoRA~\cite{hu2022lora} adapts a frozen model through low-rank updates, training a
small fraction of the parameters. It is normally motivated by cost. Our results
suggest a second reason to prefer it here: with 1.13\% of parameters trainable, the
adapted model gains code-mixed performance \emph{and keeps} English, whereas
fine-tuning the whole encoder on identical data reaches a better seen-tool error and
loses English almost entirely (Section~\ref{sec:systems}). Restricting capacity
appears to limit catastrophic forgetting in a way that matters for deployment.

\subsection{The attacks}

We generate attacks rather than reuse a fixed attack set, because the code-mixed
attacks we need do not exist as a benchmark. XTTS-v2~\cite{casanova2024xtts} provides
zero-shot cross-lingual cloning from a few seconds of reference audio;
RVC~\cite{rvc2023} provides voice conversion, which changes speaker identity while
preserving the original delivery and therefore probes a different cue;
Tortoise-TTS~\cite{betker2023tortoise} is held out entirely and used only as an
unseen tool at evaluation. AffectDF~\cite{affectdf2026} informs our attack taxonomy
and reporting format.

\subsection{Code-mixed speech resources}

MUCS 2021~\cite{diwan2021mucs} supplies the Hindi--English speech this work is built
on. Code-mixed corpora are scarce relative to monolingual ones, which is itself part
of the problem this paper describes: the detector cannot be evaluated on a language
setting for which no evaluation data has been assembled.

\subsection{What this survey does not cover}

This is a targeted review of the work our system is built from, not a complete
survey of audio deepfake detection. In particular we do not review the wider
countermeasure literature beyond the self-supervised line, nor the speaker
verification literature that anti-spoofing was originally developed to protect, nor
detection for languages other than English and Hindi--English. We also found no
prior work measuring detector generalisation to code-mixed speech over a telephone
channel; we state that as the gap we address rather than as a claim of novelty we
can prove exhaustively.
